\documentclass[a4paper]{report}
\usepackage{a4wide}
\usepackage[utf8]{inputenc}
\usepackage{parskip}
\usepackage{hyperref}
\usepackage{epsfig}
\usepackage{background}
\usepackage{mathptmx}

% To avoid tikz error, see https://tex.stackexchange.com/questions/165929/semiverbatim-with-tikz-in-beamer
\makeatletter
\global\let\tikz@ensure@dollar@catcode=\relax
\makeatother

\backgroundsetup{
scale=1,
angle=0,
opacity=1,
contents={\includegraphics[width=\paperwidth,height=\paperheight]{images/spi-front.jpg}}
}

\hypersetup{
  colorlinks   = true,
  urlcolor     = blue,
  linkcolor    = blue,
  pdfinfo = {
    Title = {SPI Annual Report 2021},
    Author = {Software in the Public Interest, Inc.},
    Keywords = {SPI, free software, open source, FOSS, annual report, charity, non-profit, 501c3},
  }
}

\begin{document}

\title{Software in the Public Interest, Inc.\\
2021 Annual Report}
\date{July XXX, 2022}

\maketitle

\newpage

\backgroundsetup{
scale=1,
angle=0,
opacity=1,
contents={\includegraphics[width=\paperwidth,height=\paperheight]{images/spi-content.jpg}}
}

\hspace{1em}

To the membership, board and friends of Software in the Public Interest, Inc:

As mandated by Article 8 of the SPI Bylaws, I respectfully submit this annual report on the activities of Software in the Public Interest, Inc. and extend my thanks to all of those who contributed to the mission of SPI in the past year.

  \emph{-- Michael Schultheiss, SPI President}

\newpage

\tableofcontents

\newpage

\chapter{President's Welcome}
\label{sec:president}

The mission of SPI is to help substantial and significant open source
projects by handling their non-technical administrative tasks so that
they aren't required to operate their own legal entity.

During the current board term, SPI (its officers, directors, contractors
and volunteers) has continued to refine and improve its processes and
response time.  Due to the continuing global pandemic, we were unable to
have our customary in-person board meetings. We were able to schedule
a virtual board sprint in January 2022 and an in-person Treasury sprint
in May 2022. We hope to return to our annual in-person sprints and
meetings with a larger set of the board when possible as they continue
to be highly productive and I strongly recommend they continue going
forward when in person meetings are again practical. SPI Vice President
Stephen Frost has been instrumental in refining our contracting process
and has been an excellent liaison with SPI's legal counsel.

SPI continues to fill a needed role in the free and open source software
community. Sister organizations exist and are not competitors for
hosting projects but have differing roles in the ecosystem. We welcome
associated projects to freely move to other umbrella organizations if
they need more than SPI's relatively hands off manner can provide as
well as welcome any projects who prefer to join SPI due to our focus.

SPI had its financial statements independently audited for the 2020 calendar
year and plans to continue having its financial statements independently
audited annually going forward. The independent audit for the 2021 calendar
year is currently underway.

The SPI board is extremely thankful for the work Martin Michlmayr has
done working with SPI's auditor and streamlining our bookkeeping. Martin
has been instrumental in releasing periodic treasury reports, updating
the financial toolset and assisting the treasury team with data
validation.

I have been involved with SPI and a number of its associated projects
for many years and am pleased the SPI continues to grow and thrive.
I look forward to continued growth of SPI over the coming year.  Thank
you again to all SPI officers, directors, contractors, volunteers and
members. SPI continues to thrive thanks to your contributions.

  \emph{-- Michael Schultheiss, SPI President}

\chapter{Committee Reports}
\section{Membership Committee}

\subsection{Statistics}

On January 1, 2021 we had 232 contributing and 1173 non-contributing members.  On December 31, 2021 there were 243 contributing members and 1253 non-contributing members.  This is an increase of 11 contributing members and an increase of 80 non-contributing members.

\chapter{Board Report}
\section{Board Members}

Board members as of January 1, 2021:

\begin{itemize}
\item Michael Schultheiss (President)
\item Stephen Frost (Vice President)
\item Tim Potter (Secretary)
\item Martin Zobel-Helas (Treasurer)
\item Joe Conway
\item Luca Filipozzi
\item Forrest Fleming
\item Chris Lamb
\item Héctor Orón Martínez
\end{itemize}

Board members as of December 31, 2021:

\begin{itemize}
\item Michael Schultheiss (President)
\item Stephen Frost (Vice President)
\item Tim Potter (Secretary)
\item Héctor Orón Martínez (Treasurer)
\item Joe Conway
\item Forrest Fleming
\item Milan Kupcevic
\item Chris Lamb
\item Martin Zobel-Helas
\end{itemize}

\section{Board Changes}

Changes that occurred during the year:

\begin{itemize}

\item Luca Filipozzi resigned from the board in February 2021.  We'd like to thank Luca for his contributions!

\item The board appointed Milan Kupcevic as an interim director in April 2021.

\item On August 16, 2021 the board voted to appoint the following officers:

\begin{itemize}
\item President: Michael Schultheiss
\item Vice President: Stephen Frost
\item Secretary: Tim Potter
\end{itemize}

\item On September 13, 2021 the board voted to appoint the following officers:

\begin{itemize}
\item Treasurer: Héctor Orón Martínez
\end{itemize}

\end{itemize}

\section{Elections}

A board membership election was conducted in July 2021.  There were 3 board seats up for election.  Nominations were received from Stephen Frost, Milan Kupcevic, and Michael Schultheiss.  Since there were 3 nominations for 3 board seats, no vote was required and all three candidates were elected for a 3 year term.

\chapter{Treasurer's Report}

\chapter{Member Project Reports}

\section{New Associated Projects}

\subsection{Adélie Linux}

\href{https://www.adelielinux.org/}{Adélie} is an independent Linux distribution committed to integrity, privacy, and user freedom.  Adélie is lightweight and supports many 32- and 64-bit computer architectures. Adélie is designed for both desktop and server use.

\subsection{GeekBeacon}

FIXME

\subsection{PMIx}

The \href{https://pmix.org/}{PMIx Standard Community Project} is governed by the PMIx Administrative Steering Committee (ASC). The ASC is a formal body that manages the PMIx Standard, PMIx Governance, and associated documents. The PMIx Standard document defines an API that provides libraries, tools, and resource managers with portable and well-defined access to commonly needed services in high-performance distributed and parallel computing systems.

\section{Projects No Longer Associated with SPI}

\begin{itemize}

\item OpenWrt joined Software Freedom Conservancy.

\item Performance Co-Pilot started using the OpenCollective Foundation platform to manage funds and no longer needs services from SPI.

\end{itemize}

\section{Updates from Associated Projects}


\appendix
\chapter{About SPI}

SPI is a non-profit organization which was founded to help organizations develop and distribute open hardware and software. We encourage programmers to use the GNU General Public License or other licenses that allow free redistribution and use of software, and hardware developers to distribute documentation that will allow device drivers to be written for their product.

SPI was incorporated as a non-profit organization on June 16, 1997 in the state of New York. Since then, it has become an umbrella organization for projects from the community.

In 1999, the Internal Revenue Service (IRS) of the United States government determined that under section 501(a) of the Internal Revenue Code SPI qualifies for 501(c)(3) (non-profit organization) status under section 509(a)(1) and 170(b)(1)(A)(vi). This means that donations made to SPI and its supported projects are tax-deductible as charitable donations for US taxpayers.

\newpage

\pagestyle{empty}

\backgroundsetup{
scale=1,
angle=0,
opacity=1,
contents={\includegraphics[width=\paperwidth,height=\paperheight]{images/spi-back-2020.jpg}}
}

\null

\end{document}
% Keep this at the bottom, thanks.
% Local Variables:
% TeX-master: "report"
% End:
